% ==============================================================================
% Section 1: Introduction, System Overview & Methodological Motivation
% ==============================================================================

\section{Introduction}
\label{sec:introduction}

Sessile autotrophs must partition finite photosynthetic income among vegetative growth, reproductive dispersal, and defence \cite{gatehouse2002plant, farmer2001surface, baldwin2006volatile}. Allocating carbon to constitutive morphological barriers --- lignified cell walls, trichomes, and biogenic silica --- or to the de-novo synthesis of volatile organic compounds (VOCs) and secondary toxic metabolites lowers somatic growth rates and seed output \cite{gatehouse2002plant, babikova2013underground, arimura2025cracking}. Minimal defence investment, in contrast, leaves plants exposed to severe defoliation and mortality under sustained herbivore outbreaks \cite{holling1959components, charnov1976optimal, rosenzweig1963graphical}.

\subsection{Limitations of Classical Population Models}
\label{sec:intro_fallacy_mixing}

Mathematical biology has long modelled consumer--resource dynamics primarily through continuous-time ordinary differential equations (ODEs), following Lotka \cite{lotka1925elements}, Volterra \cite{volterra1926fluctuations}, and Rosenzweig and MacArthur \cite{rosenzweig1963graphical}. Although mean-field formulations offer analytical tractability and closed-orbit limit cycles, they assume instantaneous, spatially homogeneous mass-action mixing that collapses the biotope into a single point.

In spatial ecosystems, three physical factors invalidate homogeneous mixing \cite{skellam1951random, turing1952chemical, okubo2001diffusion}. First, sessile hosts occupy fixed coordinates, generating localised grazing clusters and spatial refugia in which vegetation recovers while herbivores cross barren ground \cite{charnov1976optimal}. Second, heterotrophic insects lack global resource awareness and navigate boundary layers via local olfactory chemotaxis \cite{keller1970initiation, okubo2001diffusion}, sampling nearby attractants and deterrent toxins. Third, inter-plant warning propagation incurs physical delays across two distinct media: airborne plumes governed by advection--diffusion--reaction partial differential equations (PDEs) \cite{farmer2001surface, aratani2023green, arimura2025cracking}, and subterranean common mycelial networks (CMNs) governed by a cardinal-adjacent dynamic root-link graph \cite{babikova2013underground, karst2023positive, henriksson2023reexamining}. Once spatial lag, finite transit speed, and synthesis latencies become material, purely non-spatial ODE approximations lose predictive accuracy.

\begin{figure*}[t]
  \centering
  % ==============================================================================
% Figure 1: Coupled Hybrid Architecture Schematic (Native TikZ/PGF)
% ==============================================================================

\resizebox{0.96\textwidth}{!}{\begin{tikzpicture}[
    >=Stealth,
    font=\scriptsize,
    card/.style={draw=#1!65!black, fill=#1!5, rounded corners=2pt, line width=0.65pt, inner sep=6pt, align=left, text width=5.05cm, minimum height=2.55cm},
    layer/.style={draw=black!45, dashed, fill=black!2, rounded corners=3pt, line width=0.6pt, inner xsep=10pt, inner ysep=10pt},
    badge/.style={draw=#1!65!black, fill=white, rounded corners=1.5pt, line width=0.5pt, inner sep=2pt, font=\scriptsize\bfseries, text=#1!55!black}
  ]

  % Top Layer: Discrete ECS Cards
  \node[card=teal] (flora) {%
    \centering\textbf{\footnotesize\color{teal!55!black}Autotrophic Flora Entities}\\[4pt]
    \raggedright
    $\bullet$ Contiguous 1D/2D SoA buffers\\
    $\bullet$ Reserve $E_i(t) \in [0, E_{\max}]$\\
    $\bullet$ Armour $\delta_{\mathrm{morph}} \in [0, 0.50]$\\
    $\bullet$ Delay $\theta_{\mathrm{synth}} \in [8, 12]$ ticks%
  };

  \node[card=orange, right=1.85cm of flora] (def) {%
    \centering\textbf{\footnotesize\color{orange!65!black}Biotic Defence Engine}\\[4pt]
    \raggedright
    $\bullet$ Holling Type~II \& wear $c_{\mathrm{bite}}$\\
    $\bullet$ Elicitor trigger $\to$ alert state\\
    $\bullet$ Primed: $\theta_{\mathrm{primed}} \ll \theta_{\mathrm{synth}}$\\
    $\bullet$ Symbiotic tax $\delta_E^{\mathrm{tax}}$ allocation%
  };

  \node[card=purple, right=1.85cm of def] (herb) {%
    \centering\textbf{\footnotesize\color{purple!65!black}Heterotrophic Swarms}\\[4pt]
    \raggedright
    $\bullet$ Kinematic trajectory arrays\\
    $\bullet$ Headcount $N_j(t)$, reserve $E_j(t)$\\
    $\bullet$ Displacement $\mathbf{d}_t$, period $v_j$\\
    $\bullet$ MVT-inspired patch departure%
  };

  % Bottom Layer: Continuous Eulerian Cards (ample vertical gap for conduits & Layer 2 label)
  \node[card=blue, below=2.35cm of flora] (voc) {%
    \centering\textbf{\footnotesize\color{blue!60!black}Atmospheric VOC Dispersion}\\[4pt]
    \raggedright
    $\bullet$ Wind advection $\mathbf{u}_{\mathrm{wind}}$\\
    $\bullet$ $5 \times 5$ Gaussian stencil $\mathcal{K}_{\mathrm{iso}}$\\
    $\bullet$ Photolytic clearance $\lambda_s$\\
    $\bullet$ Downwind conspecific plumes%
  };

  \node[card=cyan, below=2.35cm of def] (myco) {%
    \centering\textbf{\footnotesize\color{cyan!55!black}Vascular \& Hyphal Conduits}\\[4pt]
    \raggedright
    $\bullet$ Phloem shift $\eta(t) \to \eta_{\min}$\\
    $\bullet$ Hyphal graph $\mathbf{A}_{uv}$, gain $g_{\mathrm{relay}}$\\
    $\bullet$ Direct unattenuated warning\\
    $\bullet$ Symbiotic link tax $\delta_E^{\mathrm{tax}}$%
  };

  \node[card=violet, below=2.35cm of herb] (chemo) {%
    \centering\textbf{\footnotesize\color{violet!65!black}Chemotactic Potential Field}\\[4pt]
    \raggedright
    $\bullet$ Surface $G_t = \alpha E_{\mathrm{plant}}\eta - \beta\sum T_k$\\
    $\bullet$ Jacobi flow $F_t$ \& camo $\Phi_{\mathrm{camo}}$\\
    $\bullet$ Softmax heading over $\mathcal{V}$\\
    $\bullet$ Dynamic sensory landscape%
  };

  % Outer layer frames via background
  \begin{scope}[on background layer]
    \node[layer, fit=(flora)(def)(herb), inner ysep=14pt] (L1) {};
    \node[layer, fit=(voc)(myco)(chemo), inner ysep=14pt] (L2) {};
  \end{scope}

  % Explicit foreground layer titles positioned clearly outside/above and below frame borders
  \node[fill=white, anchor=south west, font=\small\bfseries, text=black!85, inner xsep=4pt, inner ysep=2pt] (L1_lbl) at ([xshift=10pt, yshift=3pt]L1.north west) {Lagrangian Discrete ECS Layer $\mathcal{E}_t$ (Contiguous SoA Buffers)};
  \node[fill=white, anchor=north west, font=\small\bfseries, text=black!85, inner xsep=4pt, inner ysep=2pt] (L2_lbl) at ([xshift=10pt, yshift=-3pt]L2.south west) {Eulerian Continuous Field Layer $\mathcal{G}_t$ (Double-Buffered Torus $\Omega = \mathbb{T}^2$)};

  % Intra-layer horizontal arrows
  \draw[->, line width=0.8pt, color=purple!65!black] (herb.west) -- node[badge=purple] {Grazing} (def.east);
  \draw[->, line width=0.8pt, color=teal!65!black] (def.west) -- node[badge=teal] {Feedback} (flora.east);
  \draw[<->, line width=0.8pt, color=cyan!55!black] (voc.east) -- node[badge=cyan] {Conduits} (myco.west);
  \draw[->, line width=0.8pt, color=violet!65!black] (myco.east) -- node[badge=violet] {Guidance} (chemo.west);

  % Inter-layer vertical coupling arrows cleanly separated in the 2.35cm middle corridor
  \draw[->, dashed, line width=0.75pt, color=teal!65!black] ([xshift=-1.2cm]flora.south) -- node[badge=teal] {$Q_s$ Emit} ([xshift=-1.2cm]voc.north);
  \draw[->, dashed, line width=0.75pt, color=blue!60!black] ([xshift=1.2cm]voc.north) -- node[badge=blue] {Priming $S(C)$} ([xshift=1.2cm]flora.south);

  \draw[->, dashed, line width=0.75pt, color=orange!65!black] ([xshift=-1.2cm]def.south) -- node[badge=orange] {Phloem $\Delta\eta$} ([xshift=-1.2cm]myco.north);
  \draw[->, dashed, line width=0.75pt, color=cyan!55!black] ([xshift=1.2cm]myco.north) -- node[badge=cyan] {Hyphal Alert} ([xshift=1.2cm]def.south);

  \draw[->, dashed, line width=0.75pt, color=violet!65!black] ([xshift=-1.2cm]chemo.north) -- node[badge=violet] {Olfaction $\nabla F_t$} ([xshift=-1.2cm]herb.south);
  \draw[->, dashed, line width=0.75pt, color=purple!65!black] ([xshift=1.2cm]herb.south) -- node[badge=purple] {Depletion $\Delta E$} ([xshift=1.2cm]chemo.north);

  % Explicit vertical padding so outer north/south layer headers never clip
  \path ([yshift=10pt]current bounding box.north) ([yshift=-10pt]current bounding box.south);
\end{tikzpicture}
}

  \caption{The Coupled Hybrid Dynamical Architecture of PHIDS. Discrete biological entities stored in the data-oriented Entity--Component--System (ECS) layer $\mathcal{E}_t$ interact via localised grazing and damage triggers, sourcing continuous chemical concentration fields $C_s$ and guidance potentials $F_t$ evaluated across double-buffered cellular layers $\mathcal{G}_t$ on a toroidal topology $\mathbb{T}^2$.}
  \label{fig:architectural_paradigm}
\end{figure*}

\subsection{Coupled Hybrid Architecture}
\label{sec:intro_architectural_vision}

To combine spatial--biochemical detail with predictable execution throughput, we present \textbf{PHIDS} (Plant--Herbivore Interaction \& Defence Simulator). Extending our earlier discrete-event prototype \cite{dietrich2025visual}, PHIDS implements a coupled hybrid dynamical system \cite{swat2012multiscale, topping2003almass, bravo2020hybrid, husar2024mcell4} partitioned into two synchronised layers (Fig.~\ref{fig:architectural_paradigm}):
\begin{enumerate}
  \item \textbf{Discrete Lagrangian Entity Layer ($\mathcal{E}_t$):} Sessile plants and mobile herbivore swarms reside in a data-oriented Entity--Component--System (ECS) backed by contiguous Structure-of-Arrays (SoA) memory buffers \cite{lam2015numba, harris2020array}, avoiding per-organism pointer chasing and supporting $\mathcal{O}(1)$ spatial lookups.
  \item \textbf{Continuous Eulerian Field Layer ($\mathcal{G}_t$):} Semiochemicals, volatile concentrations, and sensory fields are stored as 2D scalar arrays on a periodic toroidal lattice $\Omega = \mathbb{T}^2$. Airborne transport evolves via operator-split Reaction--Diffusion--Advection PDEs \cite{strang1968construction, leveque2002finite}, whereas chemotactic navigation potentials are computed via compiled Jacobi iterative relaxation kernels \cite{lam2015numba}.
\end{enumerate}

By coupling discrete agent transitions with continuous field operators at hourly resolution, PHIDS resolves Holling Type~II intake saturation \cite{holling1959components}, an MVT-inspired local patch-departure heuristic \cite{charnov1976optimal}, sigmoidal Hill receptor priming, and vascular phloem withdrawal \cite{schwachtje2006snf1} across tens of thousands of interacting organisms.

\subsection{Structural Homomorphism to Distributed WSN Security}
\label{sec:intro_wsn_paradigm}

Beyond spatial systems biology, the dynamical architecture of PHIDS shares a formal structural homomorphism and computational Design Space Exploration (DSE) methodology with distributed resilience planning in large-scale static Wireless Sensor Networks (WSNs) \cite{forster2023novel, forster2024determining, forster2024topology, dressler2010survey}. Static sensor deployments face a parallel energy-allocation trade-off to sessile flora (Table~\ref{tab:energetic_allocation_tradeoffs}): continuous cryptographic checks deplete constrained battery reserves, whereas passive sleep states expose networks to multi-hop disruption during activation delays \cite{forster2023novel, forster2024determining}. Building on our prior WSN security and EEDSE work \cite{forster2023novel, forster2024determining, forster2024topology, dietrich2025visual}, Section~\ref{sec:wsn_isomorphism} formalises this structural mapping --- including explicit non-isomorphism boundaries --- and transfers the multi-fidelity calibration, sensitivity screening, and deterministic SoA workflows across both domains.

\begin{table*}[b]
  \centering
  \caption{Energetic Allocation Trade-offs Across Autotrophic Phenotypic Strategies in Meso-Scale Biotopes.}
  \label{tab:energetic_allocation_tradeoffs}
  % ==============================================================================
% Table 1: Energetic Allocation Trade-offs Across Autotrophic Phenotypes
% Placed at [b] (bottom of page)
% ==============================================================================

\footnotesize
\renewcommand{\arraystretch}{1.15}
\begin{tabularx}{\textwidth}{
  >{\raggedright\arraybackslash\hsize=0.85\hsize}X
  >{\centering\arraybackslash\hsize=0.70\hsize}X
  >{\centering\arraybackslash\hsize=0.68\hsize}X
  >{\centering\arraybackslash\hsize=0.68\hsize}X
  >{\raggedright\arraybackslash\hsize=0.74\hsize}X
  >{\raggedright\arraybackslash\hsize=0.85\hsize}X
  >{\raggedright\arraybackslash\hsize=1.50\hsize}X
}
  \toprule
  \textbf{Plant Phenotype Class} &
  \textbf{Vegetative Growth} ($g_j$) &
  \textbf{Morph. Upkeep} ($\delta_{\text{morph}}$) &
  \textbf{Inducible Defence} ($\delta_{\text{chem}}$) &
  \textbf{Signalling Channel} &
  \textbf{Target Herbivore Guild} &
  \textbf{Ecological Strategy} \\
  \midrule
  \textbf{Pioneer Autotroph} &
  High ($85$--$95\%$) &
  Negligible ($0$--$2\%$) &
  Zero ($0\%$) &
  None &
  Generalist Foliivores &
  Rapid spatial colonisation prioritising reproductive output over individual tissue protection. \\
  \addlinespace[3pt]
  \textbf{Constitutive Armoured} &
  Low ($20$--$35\%$) &
  High ($30$--$50\%$) &
  Low ($5$--$10\%$) &
  Passive Airborne &
  Chewing Macrobrowsers &
  Static physical hardening paying sustained maintenance costs for zero-latency protection. \\
  \addlinespace[3pt]
  \textbf{Inducible Synthesiser} &
  Medium ($50$--$65\%$) &
  Low ($5$--$10\%$) &
  High ($25$--$40\%$) &
  Airborne VOC Plumes &
  Specialist Swarms &
  On-demand allelochemical synthesis, vulnerable during countdown latency $\theta_{\text{synth}}$ unless primed. \\
  \addlinespace[3pt]
  \textbf{Symbiotic Networker} &
  Medium ($45$--$60\%$) &
  Low ($5$--$10\%$) &
  Medium ($20$--$30\%$) &
  Subterranean Hyphae &
  Invasive Phloem Feeders &
  Maintains persistent CMN conduits at carbon tax $\delta_E^{\text{tax}}$ to sustain wind-independent local warnings. \\
  \bottomrule
\end{tabularx}

\end{table*}

\subsection{Summary of Contributions}
\label{sec:intro_contributions}

The contributions of this work are fourfold:
\begin{enumerate}
  \item \textbf{Coupled Hybrid Simulator Architecture:} A data-oriented SoA Entity--Component--System paired with operator-split 2D reaction--diffusion--advection PDEs, Jacobi flow relaxation, and dynamic cardinal root graphs, compiled to LLVM native kernels with bit-exact thread determinism and Zarr compression (Section~\ref{sec:computational_realisation}).
  \item \textbf{Multi-Channel Botanical Defence Kinetics:} Parameterised operational equations for constitutive armour, structural tissue upkeep, airborne volatile dispersion, Hill receptor priming, subterranean mycorrhizal relays, and phloem palatability withdrawal.
  \item \textbf{Analytical System Bounds:} Three analytical propositions establishing carrying-capacity boundedness and non-reproductive dissipation, the closed-form continuous-grazing critical synthesis delay $\theta_{\text{crit}}$, and a mean-field delay-bifurcation criterion ($\tau^* = 2.38\,\text{h} > 1.0\,\text{h}$), checked against empirical cross-thread invariance.
  \item \textbf{In-Silico Case Studies and Structural DSE Mapping:} Evaluation on a $256 \times 256$ Central European ecotone spanning multi-fidelity Bayesian calibration, $10$-seed single-season zero-shot transfer, continuous-vs-diurnal bound diagnostics, $90$-run dual-channel wind ablations ($+21$-run upwind relay), and $1{,}792 + 224$-run CRN Sobol sensitivity analyses alongside formal structural mapping to static WSN DSE pipelines.
\end{enumerate}

\subsection{Article Organisation}
\label{sec:intro_organization}

Section~\ref{sec:biological_domain} introduces the biological defence tiers and foraging kinematics, and Section~\ref{sec:related_works} positions PHIDS relative to existing simulation frameworks. Section~\ref{sec:mathematical_framework} formulates the coupled hybrid mathematical model, Section~\ref{sec:theoretical_invariants} derives the analytical bounds and stability properties, and Section~\ref{sec:case_studies} evaluates the framework on the Central European ecotone. Section~\ref{sec:computational_realisation} details the ECS memory layout, JIT kernel design, Roofline scaling, double-buffered determinism, and Zarr replay storage; Section~\ref{sec:discussion} synthesises ecological buffering, the structural WSN homomorphism and DSE transfer, and model assumptions; and Section~\ref{sec:conclusion} concludes.
